\setcounter{section}{16}






  \zwischenueberschrift{Übungsaufgaben}

\inputaufgabegibtloesung
{}
{





Es sei

\mavergleichskette
{\vergleichskette
{ X
}
{ = }{ K\!-\!\operatorname{Spek}\, \left(  R  \right)
}
{  }{
}
{    }{
}
{   }{
}
}
{}{}{}
das
$K$-\definitionsverweis {Spektrum}{}{}
einer
\definitionsverweis {endlich erzeugten}{}{}
kommutativen $K$-Algebra. Zeige, dass ein
\definitionsverweis {irreduzibler Filter}{}{}
durch
\definitionsverweis {offene Mengen}{}{}
der Form $D(f)$ erzeugt wird.





}
{} {}

\inputaufgabe
{}
{





Es sei

\mavergleichskette
{\vergleichskette
{ X
}
{ = }{ K\!-\!\operatorname{Spek}\, \left(  R  \right)
}
{  }{
}
{    }{
}
{   }{
}
}
{}{}{}
eine
\definitionsverweis {affine Varietät}{}{}
und

\mavergleichskette
{\vergleichskette
{ Z
}
{ \subseteq }{ X
}
{  }{
}
{    }{
}
{   }{
}
}
{}{}{}
eine
\definitionsverweis {abgeschlossene}{}{}
Teilmenge. Zeige, dass der
\definitionsverweis {Umgebungsfilter}{}{}
$U(Z)$ von
\definitionsverweis {offenen Mengen}{}{}
der Form $D(f)$ erzeugt wird.





}
{} {}

\inputaufgabe
{}
{





Es sei $X$ ein topologischer Raum. Zeige, dass ein
\definitionsverweis {Ultrafilter}{}{}
\definitionsverweis {irreduzibel}{}{}
ist.





}
{} {}

\inputaufgabe
{}
{





Zeige, dass sich in der Korrespondenz aus

Satz 16.2

\definitionsverweis {maximale Ideale}{}{,}
Punkte im
$K$-\definitionsverweis {Spektrum}{}{}
und
\definitionsverweis {Umgebungsfilter}{}{}
von Punkten entsprechen.





}
{} {}

\inputaufgabe
{}
{





Zeige, dass sich in der Korrespondenz aus

Satz 16.2

\definitionsverweis {minimale Primideale}{}{,}

\definitionsverweis {irreduzible Komponenten}{}{}
des
$K$-\definitionsverweis {Spektrums}{}{}
und
\definitionsverweis {Ultrafilter}{}{}
entsprechen.





}
{} {}

\inputaufgabe
{}
{





Es sei $K$ ein
\definitionsverweis {algebraisch abgeschlossener Körper}{}{,}
wir betrachten die affine Ebene ${\mathbb A}^{2}_{K}$ zusammen mit der $x$-Achse

\mavergleichskettedisp
{\vergleichskette
{ V
}
{ =} { V(y)
}
{ } {
}
{ } {
}
{ } {
}
}
{}{}{.}
Zeige, dass die folgende Menge ein saturiertes multiplikatives System ist.

\mavergleichskettedisp
{\vergleichskette
{S
}
{ =} { \left\{  f \in K[X,Y]   \mid   \text{In der homogenen Komponente } f_{\deg(f)} \text{ kommt } x^{\deg(f)} \text{ vor}         \right\}
}
{ } {
}
{ } {
}
{ } {
}
}
{}{}{.}
Skizziere die Nullstellenmenge von einigen Polynomen, die oder die nicht zu $S$ gehören.

Es sei $F$ der zugehörige topologische Filter. Vergleiche $F$ mit dem Umgebungsfilter zu $V$ und dem generischen Filter zu $V$.





}
{} {}

\inputaufgabe
{}
{





Es sei $U$ eine
\definitionsverweis {quasiaffine Varietät}{}{}
über einem
\definitionsverweis {algebraisch abgeschlossenen Körper}{}{}
$K$. Zeige, dass die
\definitionsverweis {Einheiten}{}{}
in $\Gamma (U, {\mathcal O}_U )$ den
\definitionsverweis {Morphismen}{}{}
von $U$ nach

\mavergleichskette
{\vergleichskette
{ { {\mathbb A}_{ K }^{ \times  } }
}
{ = }{ {\mathbb A}^{1}_{K} \setminus \{0\}
}
{  }{
}
{    }{
}
{   }{
}
}
{}{}{}
entsprechen.





}
{} {}

\inputaufgabe
{}
{





Es sei $U$ eine
\definitionsverweis {quasiaffine Varietät}{}{}
über einem
\definitionsverweis {algebraisch abgeschlossenen Körper}{}{}
$K$ und sei
\maabbdisp {{\psi}} { U } { { {\mathbb A}_{  K }^{  n   } }
} {}
ein
\definitionsverweis {Morphismus}{}{.}
Zeige, dass ${\psi}$ genau dann durch die abgeschlossene Menge

\mavergleichskette
{\vergleichskette
{ V({\mathfrak a})
}
{ \subseteq }{ { {\mathbb A}_{  K }^{  n   } }
}
{  }{
}
{    }{
}
{   }{
}
}
{}{}{}
faktorisiert, wenn ${\mathfrak a}$ im
\definitionsverweis {Kern}{}{}
des globalen Ringhomomorphismus
\maabbdisp {\tilde{\psi}} { K [ T_1, \ldots ,  T_{ n }]  } { \Gamma (U, {\mathcal O}_U )
} {}
liegt.





}
{} {}

\inputaufgabe
{}
{





Es seien $U$ und $V$
\definitionsverweis {quasiaffine Varietäten}{}{}
über einem
\definitionsverweis {algebraisch abgeschlossenen Körper}{}{}
$K$ und sei

\mavergleichskette
{\vergleichskette
{W
}
{ = }{ U \uplus V
}
{  }{
}
{    }{
}
{   }{
}
}
{}{}{}
ihre
\definitionsverweis {disjunkte Vereinigung}{}{.}
Zeige, dass für jede quasiaffine Varietät $Z$ ein
\definitionsverweis {Morphismus}{}{}
von $W$ nach $Z$ einfach ein Paar bestehend aus einem Morphismus von $U$ nach $Z$ und einem Morphismus von $V$ nach $Z$ ist.





}
{} {}

\inputaufgabegibtloesung
{}
{





Man gebe ein Beispiel von zwei
\definitionsverweis {affinen Varietäten}{}{}
$V_1$ und $V_2$ und einem
\definitionsverweis {Morphismus}{}{}
\maabbdisp {{\psi}} {V_1 } { V_2
} {,}
der bijektiv ist, wo aber die Umkehrabbildung nicht stetig  ist.





}
{} {}

\inputaufgabegibtloesung
{}
{





Es sei

\mavergleichskette
{\vergleichskette
{ V
}
{ = }{ V \left(  x^2+y^2-1  \right)
}
{ \subseteq }{ {\mathbb A}^{2}_{ K }
}
{    }{
}
{   }{
}
}
{}{}{}
der Einheitskreis über einem Körper $K$ und es seien
\mathkor {} {P=(a,b)} {und} {Q=(c,d)} {}
Punkte auf $V$. Zeige, dass es einen
\definitionsverweis {Automorphismus}{}{}
\maabb {\varphi} { V } { V
} {}
mit

\mavergleichskettedisp
{\vergleichskette
{ \varphi(P)
}
{ =} { Q
}
{ } {
}
{ } {
}
{ } {
}
}
{}{}{}
gibt.





}
{} {}

\inputaufgabegibtloesung
{}
{





\aufzaehlungdrei{Skizziere die Nullstellengebilde

\mavergleichskettedisp
{\vergleichskette
{V
}
{ =} { V(XY,XZ,YZ)
}
{ \subseteq} { { {\mathbb A}_{ K }^{  3   } }
}
{ } {
}
{ } {
}
}
{}{}{}
und

\mavergleichskettedisp
{\vergleichskette
{W
}
{ =} { V(ST(S-T))
}
{ \subseteq} { {\mathbb A}^{2}_{K}
}
{ } {
}
{ } {
}
}
{}{}{}
im reellen Fall.
}{Stifte einen bijektiven Morphismus
\maabbdisp {\varphi} { V } { W
} {.}
}{Zeige, dass der Morphismus $\varphi$ außerhalb des Nullpunktes ein Isomorphismus ist
\zusatzklammer {die Charakteristik des Körpers sei $\neq 2$} {} {.}
}





}
{} {}

\inputaufgabegibtloesung
{}
{





Zeige, dass durch
\maabbeledisp {\varphi} { V \left(  Z^2+W^2-1  \right)  } { V \left(  X^2+Y^2-1  \right)
} { (Z,W) } {  \left(  Z^2 - W^2   , \,   2ZW                   \right)   = (X,Y)
} {.}
ein
\definitionsverweis {Morphismus}{}{}
des Einheitskreises in sich gegeben ist. Zeige, dass das Urbild zu jedem Punkt

\mavergleichskette
{\vergleichskette
{ P
}
{ \in }{ V \left(  X^2+Y^2-1  \right)
}
{  }{
}
{    }{
}
{   }{
}
}
{}{}{}
aus zwei Punkten besteht.





}
{} {}

\inputaufgabe
{}
{





Es seien $X$ und $Z$
\definitionsverweis {quasiaffine Varietäten}{}{}
über einem
\definitionsverweis {algebraisch abgeschlossenen Körper}{}{}
$K$. Zu jeder offenen Teilmenge

\mavergleichskette
{\vergleichskette
{U
}
{ \subseteq }{ X
}
{  }{
}
{    }{
}
{   }{
}
}
{}{}{}
betrachten wir

\mavergleichskettedisp
{\vergleichskette
{ { \mathcal  F   }  \left(   U   \right)
}
{ \defeq} { \left\{  \varphi:U \rightarrow Z  \mid   \varphi \text{ ist ein Morphismus}         \right\}
}
{ } {
}
{ } {
}
{ } {
}
}
{}{}{.}
Zeige, dass dies eine
\definitionsverweis {Garbe}{}{}
auf $X$ ist.





}
{} {}

\inputaufgabegibtloesung
{}
{





Man beschreibe einen
$K$-\definitionsverweis {Algebrahomomorphismus}{}{}
derart, dass die induzierte
\definitionsverweis {Spektrumsabbildung}{}{}
der
$K$-\definitionsverweis {Spektren}{}{}
die Addition auf $K$ beschreibt.





}
{} {}

\inputaufgabe
{}
{





Es sei $K$ ein algebraisch abgeschlossener Körper. Zeige, dass die Addition, die Multiplikation, das Negative, das Inverse und die Division in $K$ sich als Morphismen realisieren lassen.





}
{} {}

\inputaufgabe
{}
{





Es sei $R$ ein
\definitionsverweis {kommutativer Ring}{}{}
und sei

\mavergleichskette
{\vergleichskette
{ {\mathfrak a}
}
{ = }{  \left(  f_1 , \ldots ,  f_n  \right)
}
{  }{
}
{    }{
}
{   }{
}
}
{}{}{}
ein endlich erzeugtes
\definitionsverweis {Ideal}{}{.}
Es sei

\mavergleichskette
{\vergleichskette
{ f
}
{ \in }{ R
}
{  }{
}
{    }{
}
{   }{
}
}
{}{}{}
ein weiteres Element. Dann nennt man die $R$-Algebra

\mavergleichskettedisp
{\vergleichskette
{A
}
{ =} { R [ T_1, \ldots ,  T_{ n }]/ \left(   f _1  T_{1}   + \cdots +   f _{ n } T_{ n }+ f   \right)
}
{ } {
}
{ } {
}
{ } {
}
}
{}{}{}
die \definitionswort {erzwingende Algebra}{} zu den
\mathl{f_1 , \ldots , f_n,f}{.} Zeige, dass $A$ folgende Eigenschaft erfüllt: Zu jedem
\definitionsverweis {Ringhomomorphismus}{}{}
\maabb {\varphi} { R } { S
} {}
in einen kommutativen Ring $S$ mit der Eigenschaft

\mavergleichskette
{\vergleichskette
{ \varphi(f)
}
{ \in }{ {\mathfrak a} S
}
{  }{
}
{    }{
}
{   }{
}
}
{}{}{}
gibt es einen
$R$-\definitionsverweis {Algebrahomo\-morphismus}{}{}
\maabb {\vartheta} { A } { S
} {.}
Zeige ebenso, dass dieser Homomorphismus
\betonung{nicht}{} eindeutig bestimmt ist.





}
{} {}

\inputaufgabe
{}
{





Es sei $R$ eine kommutative $K$-Algebra von endlichem Typ über einem algebraisch abgeschlossenen Körper. Es seien
\mathl{f_1 , \ldots ,  f_n,f}{} Elemente in $R$ und es sei

\mavergleichskettedisp
{\vergleichskette
{A
}
{ =} { R [ T_1, \ldots ,  T_{ n }]/ \left(   f _1  T_{1}   + \cdots +   f _{ n } T_{ n }+ f   \right)
}
{ } {
}
{ } {
}
{ } {
}
}
{}{}{}
die
\definitionsverweis {erzwingende Algebra}{}{}
zu diesen Daten. Charakterisiere die
\definitionsverweis {Fasern}{}{}
des zugehörigen Morphismus
\maabbdisp {} { K\!-\!\operatorname{Spek}\, \left(  A  \right) } { K\!-\!\operatorname{Spek}\, \left(  R  \right)
} {.}





}
{} {}





  \zwischenueberschrift{Aufgaben zum Abgeben}

\inputaufgabe
{4}
{





Es sei $K$ ein
\definitionsverweis {algebraisch abgeschlossener Körper}{}{}
und seien $R$ und $S$
\definitionsverweis {integre}{}{}
$K$-\definitionsverweis {Algebren von endlichem Typ}{}{.}
Es sei
\maabb {\varphi} { R } { S
} {}
ein
$K$-\definitionsverweis {Algebrahomomorphis\-mus}{}{}
mit zugehörigem Morphismus
\maabb {\varphi^*} { K\!-\!\operatorname{Spek}\, \left(  S  \right) } { K\!-\!\operatorname{Spek}\, \left(  R  \right)
} {.}
Zeige, dass folgende Aussagen äquivalent sind.
\aufzaehlungdreiabc{ $\varphi$ ist injektiv.
}{Das Bild von $\varphi^*$ ist
\definitionsverweis {dicht}{}{}
in
\mathl{K\!-\!\operatorname{Spek}\, \left(  R  \right)}{.}
}{ $\varphi$ induziert einen Ringhomomorphismus
\maabb {} { Q(R) } { Q(S)
} {.}
}





}
{} {}

\inputaufgabe
{4}
{





Es sei $K$ ein
\definitionsverweis {algebraisch abgeschlossener Körper}{}{}
und seien $R$ und $S$ zwei
\definitionsverweis {integre}{}{}
$K$-\definitionsverweis {Algebren von endlichem Typ}{}{.}
Es sei ein
$K$-\definitionsverweis {Algebrahomomorphismus}{}{}
\maabbdisp {\varphi} {Q(R) } { Q(S)
} {}
zwischen den
\definitionsverweis {Quotientenkörpern}{}{}
gegeben. Zeige, dass es eine offene Teilmenge

\mavergleichskette
{\vergleichskette
{ U
}
{ \subseteq }{ K\!-\!\operatorname{Spek}\, \left(  S  \right)
}
{  }{
}
{    }{
}
{   }{
}
}
{}{}{}
und einen Morphismus
\maabbdisp {} { U } { K\!-\!\operatorname{Spek}\, \left(  R  \right)
} {}
gibt, der $\varphi$ induziert.





}
{} {}

\inputaufgabe
{5}
{





Man gebe ein Beispiel von zwei affin-algebraischen Kurven $C_1$ und $C_2$ über $\mathbb C$ und einem \definitionsverweis {Morphismus}{}{}
\maabbdisp {{\psi}} { C_1 } { C_2
} {,}
der
\definitionsverweis {bijektiv}{}{}
ist, wo aber die Umkehrabbildung nicht stetig in der metrischen Topologie ist.





}
{} {}

\inputaufgabe
{6}
{





Man gebe ein Beispiel von zwei
\definitionsverweis {irreduziblen}{}{}
affinen Varietäten $V_1$ und $V_2$ und einem
\definitionsverweis {Morphismus}{}{}
\maabbdisp {{\psi}} {V_1 } { V_2
} {,}
der bijektiv ist, wo aber die Umkehrabbildung nicht stetig (in der Zariski-Topologie) ist.





}
{(und daher auch kein Morphismus ist).} {}

\inputaufgabe
{3}
{





Es sei

\mavergleichskette
{\vergleichskette
{ U
}
{ \subseteq }{ K\!-\!\operatorname{Spek}\, \left(   R   \right)
}
{  }{
}
{    }{
}
{   }{
}
}
{}{}{}
eine
\definitionsverweis {quasiaffine Varietät}{}{}
und sei

\mavergleichskette
{\vergleichskette
{ f
}
{ \in }{ \Gamma ( U , {\mathcal O} )
}
{  }{
}
{    }{
}
{   }{
}
}
{}{}{}
eine
\definitionsverweis {algebraische Funktion}{}{.}
Es seien

\mavergleichskette
{\vergleichskette
{ q
}
{ = }{ g_i/h_i
}
{  }{
}
{    }{
}
{   }{
}
}
{}{}{,}

\mavergleichskette
{\vergleichskette
{ i
}
{ = }{ 1 , \ldots ,  n
}
{  }{
}
{    }{
}
{   }{}
}
{}{}{,}
lokale Darstellungen von $f$ auf

\mavergleichskette
{\vergleichskette
{ D \left(  h_i  \right)
}
{ \subseteq }{ U
}
{  }{
}
{    }{
}
{   }{
}
}
{}{}{}
\zusatzklammer {die $U$ überdecken} {} {.}
Zeige, dass das Urbild $f^{-1}(0)$ gleich der abgeschlossenen Menge
\mathl{V \left(  h_1g_1 , \ldots , h_ng_n  \right) \cap U}{} ist.





}
{} {}
